\subsection{Expected Failure Modes}

File operations are kernel requests that may be refused after the Python call begins.

\subsubsection{Missing Files and \texttt{FileNotFoundError}}

Read mode on a nonexistent path raises \texttt{FileNotFoundError} with \texttt{filename}, \texttt{errno}, and \texttt{strerror}---not \texttt{None} or an empty string.

\begin{verbatim}
from pathlib import Path
path = Path("missing.txt")
try:
    path.read_text(encoding="utf-8")
except FileNotFoundError as err:
    print(err.filename, err.errno)
\end{verbatim}

\subsubsection{Permission Failures and \texttt{PermissionError}}

Existing paths may be unreadable or unwritable under current credentials. Opening a directory as a file often yields \texttt{IsADirectoryError} (an \texttt{OSError} subclass).

\subsubsection{Invalid Paths, Locked Files, and Platform-Specific Filesystem Constraints}

Embedded null bytes fail before syscall (\texttt{ValueError}). Name length limits, reserved names, and cross-process locks surface as \texttt{OSError} variants with platform-specific \texttt{errno}.
